\section{Original-record lower bounds}\label{sec:original-record-lower-bounds}

The lower bounds use a \leanref{S-4}{marked-component lower-bound experiment} within the primitive class of \cref{obj:def:primitive-class}. Under each hypothesis, a finite prior on primitive laws induces a mixture of the original sampling experiment: each primitive law generates exactly \(n\) labeled records and \(m\) outcome-free treatment records before mixing. Write \leanref{sym:Mtheta}{\ensuremath{M_\theta}} for this fixed-size mixture law under hypothesis \leanref{sym:theta}{\ensuremath{\theta\in\{-1,+1\}}}. The construction separates the point contrasts while controlling the statistical distance between the two complete-data mixtures. We use the following normalization of that distance.

\begin{definitionv}{synth_4}[Treatment-record marginal]
For a primitive Borel probability law
\(P=\mathcal L(X,A,Y(0),Y(1))\) on
\([0,1]^d\times\{0,1\}^3\), define its treatment-record marginal by
\[
P_{XA}=\mathcal L_P(X,A).
\]
Thus \(P_{XA}\) is the probability law on \([0,1]^d\times\{0,1\}\) satisfying
\[
P_{XA}(E)=P\bigl((X,A)\in E\bigr)
\]
for every Borel set \(E\subseteq[0,1]^d\times\{0,1\}\).
\end{definitionv}

\begin{definitionv}{synth_6}[Fixed-size hypothesis mixtures]
For \(n\in\mathbb N\) with \(n\ge2\), \(m\in\mathbb N_0\), and the finite priors \(\varpi_\theta\) produced by \cref{obj:def:marked-handle} on its probability-valued domain, define
\[
M_\theta
:=\int\bigl(P_o^{\otimes n}\otimes P_{XA}^{\otimes m}\bigr)\,d\varpi_\theta(P),
\qquad \theta\in\{-1,+1\},
\]
where \(P_o=\mathcal L_P(X,A,Y)\), \(Y=(1-A)Y(0)+AY(1)\), and \(P_{XA}=\mathcal L_P(X,A)\). Thus \(M_{+1}\) and \(M_{-1}\) are the positive- and negative-hypothesis probability laws of the original-record dataset \(\mathcal D_{n,m}\), with \(n\) labeled records and \(m\) outcome-free treatment records. Under each mixture, one primitive law \(P\) is drawn from the corresponding prior and all \(n+m\) records are drawn conditionally independently from that law through their respective observation channels.
\end{definitionv}

The unhalved squared Hellinger distance \leanref{sym:Htwo}{\ensuremath{H^2(M_{+1},M_{-1})}} measures the distinguishability of the two hypotheses using both record types. The finite-prior construction \leanref{sym:lowerHandle}{\ensuremath{\mathfrak C}} of \cref{obj:def:marked-handle} below couples propensity and response perturbations through correlated sign pairs. A smooth localization bump controls the contrast near the evaluation profile, while a finer smooth frame controls the spatial dependence of the perturbations.

\begin{definitionv}{def:hellinger}[Squared Hellinger distance \(H^2\)]
The squared Hellinger distance between the positive- and negative-hypothesis mixture laws \(M_{+1}\) and \(M_{-1}\) is
\[
H^2(M_{+1},M_{-1})
=
\int
\left(
\sqrt{\frac{dM_{+1}}{d(M_{+1}+M_{-1})}}
-
\sqrt{\frac{dM_{-1}}{d(M_{+1}+M_{-1})}}
\right)^2
\,d(M_{+1}+M_{-1}).
\]
\end{definitionv}

\begin{definitionv}{synth_2}[Constructed arm-mean functions]
The real-valued arm-mean functions \(\mu_{0,\theta,\eta}\) and \(\mu_{1,\theta,\eta}\) are defined for every dimension \(d\in\mathbb N\), real scales \(h,\delta\), real propensity perturbation amplitude \(a\), real response perturbation amplitude \(b\), hypothesis sign \(\theta\in\{-1,+1\}\), and response sign array \(\eta\in\{-1,+1\}^{\mathcal I}\) by
\[
\begin{aligned}
\mu_{0,\theta,\eta}(x)
&=\frac12+b\sum_{\mathbf z\in\mathcal I}
  \eta_{\mathbf z}\phi_{\mathbf z}(x)
  +\theta ab\,B_h(x)^2,\\
\mu_{1,\theta,\eta}(x)
&=\frac12+b\sum_{\mathbf z\in\mathcal I}
  \eta_{\mathbf z}\phi_{\mathbf z}(x)
  -\theta ab\,B_h(x)^2,
\qquad x\in\mathbb R^d.
\end{aligned}
\]
Here \(\mathcal I\), the finite frame index set, \(B_h\), the localization bump, and \(\phi_{\mathbf z}\), the localized frame entries, use the defining formulas in \cref{obj:def:marked-handle}, evaluated at dimension \(d\) and the stated real parameter values. The hypothesis sign and response signs encode Boolean values by assigning \(+1\) to true and \(-1\) to false. Each \(\eta_{\mathbf z}\) is the encoded second component of the Boolean sign pair at \(\mathbf z\); the first component is arbitrary. Control and treatment correspond to false and true arm values, respectively.
\end{definitionv}



In \cref{obj:def:marked-handle}, \leanref{sym:a}{\ensuremath{a}} is the propensity perturbation amplitude and \(b\) is the scalar response-mean perturbation amplitude. The arm means are centered at the shared function \(1/2+bG_\eta(x)\), with offsets \(-\tau_\theta(x)/2\) under control and \(+\tau_\theta(x)/2\) under treatment. Consequently, the shared perturbation cancels from their difference, leaving exactly \(\tau_\theta(x)\). Here the scalar \(b\) is distinct from the estimator's fine polynomial basis vector \(\mathbf b_J(x)\) in \cref{obj:def:projection}.

In \cref{obj:def:marked-handle}, the product sign-pair weights \leanref{sym:Pi}{\ensuremath{w_\theta(\lambda,\eta)}} encode the hypothesis through the correlation between the two perturbations. The construction itself depends only on \(d\) and \((h,\delta,a,b)\); the parameters \(\alpha,\beta,\gamma,L,\varepsilon\) enter only through the admissibility, occupancy, and information constants of the construction certificate. The deterministic contrast \leanref{sym:tautheta}{\ensuremath{\tau_\theta}} is shared by every primitive law within a hypothesis. Thus mixing over nuisance functions preserves an exact target separation. The compensating terms in the arm means link this separation to the singleton cancellation stated in the construction certificate below.

The information argument uses conditional factorization over spatial components. The calculus in \cref{obj:lem:hellinger-block-calculus}, stated after the construction, records how squared Hellinger distance behaves under products, conditioning on a common marginal, and the addition of a common independent coordinate.

\begin{algorithmv}{def:marked-handle}[Finite-prior construction \(\mathfrak C\)]
The construction \(\mathfrak C\) takes inputs \((h,\delta,a,b)\) satisfying \(0<\delta\le h\le1/2\) and \(a,b>0\), at the public parameters \((d,\alpha,\beta,\gamma,L)\) and fixed interior covariate profile \(x_0\).
\begin{enumerate}
\item Define the real functions
\[
\chi(u)=
\begin{cases}
\exp(-1/u),&u>0,\\
0,&u\le0,
\end{cases}
\qquad
\vartheta(u)=\frac{\pi}{2}\frac{\chi(u)}{\chi(u)+\chi(1-u)},
\qquad u\in\mathbb R,
\]
and, for each \(z\in\mathbb Z\),
\[
\omega_z(u)=
\begin{cases}
\sin\vartheta(u-z+1),&z-1\le u\le z,\\
\cos\vartheta(u-z),&z\le u\le z+1,\\
0,&u\notin[z-1,z+1].
\end{cases}
\]
The two expressions at \(u=z\) agree.

\item For \(w\in\mathbb R^d\) and \(x\in[0,1]^d\), define the localization bump and fine-scale frame by
\[
B(w)=\prod_{v=1}^d\frac{\chi(1-4w_v^2)}{\chi(1)},
\qquad
B_h(x)=B\bigl((x-x_0)/h\bigr),
\]
\[
\psi_{\mathbf z}(x)=
\prod_{v=1}^d
\omega_{z_v}\bigl((x_v-x_{0,v})/\delta\bigr),
\]
\[
\mathcal I=
\left\{\mathbf z\in\mathbb Z^d:
|z_v|\le1+\frac{h}{2\delta}
\text{ for every }v\in\{1,\ldots,d\}\right\},
\qquad
\phi_{\mathbf z}(x)=B_h(x)\psi_{\mathbf z}(x),
\quad \mathbf z\in\mathcal I.
\]

\item For sign arrays \(\lambda,\eta\in\{-1,+1\}^{\mathcal I}\), define
\[
F_\lambda(x)=
\sum_{\mathbf z\in\mathcal I}\lambda_{\mathbf z}\phi_{\mathbf z}(x),
\qquad
G_\eta(x)=
\sum_{\mathbf z\in\mathcal I}\eta_{\mathbf z}\phi_{\mathbf z}(x).
\]
For each hypothesis index \(\theta\in\{-1,+1\}\), define the sign-pair correlation and product weights by
\[
\rho_\theta=\frac{\theta}{2},
\qquad
w_\theta(\lambda,\eta)=
\prod_{\mathbf z\in\mathcal I}
\frac{1+\rho_\theta\lambda_{\mathbf z}\eta_{\mathbf z}}{4}.
\]

\item Define the propensity, contrast, and potential-response means by
\[
e_\lambda(x)=\frac12+aF_\lambda(x),
\qquad
\tau_\theta(x)=-4ab\rho_\theta B_h(x)^2,
\]
\[
\mu_{0,\theta,\eta}(x)=
\frac12+bG_\eta(x)-\frac12\tau_\theta(x),
\qquad
\mu_{1,\theta,\eta}(x)=
\frac12+bG_\eta(x)+\frac12\tau_\theta(x).
\]
The probability-valued domain consists of inputs satisfying
\[
e_\lambda(x),\ \mu_{0,\theta,\eta}(x),\
\mu_{1,\theta,\eta}(x)\in[0,1]
\quad
\text{for every }\theta\in\{-1,+1\},\
\lambda,\eta\in\{-1,+1\}^{\mathcal I},\
x\in[0,1]^d.
\]
On this domain, define the primitive Borel law \(P_{\theta,\lambda,\eta}\) by
\[
X\sim\operatorname{Unif}([0,1]^d),
\qquad
\mathcal L\bigl(A,Y(0),Y(1)\mid X=x\bigr)
=
\operatorname{Bernoulli}\bigl(e_\lambda(x)\bigr)
\otimes
\operatorname{Bernoulli}\bigl(\mu_{0,\theta,\eta}(x)\bigr)
\otimes
\operatorname{Bernoulli}\bigl(\mu_{1,\theta,\eta}(x)\bigr).
\]
Here \(X\) is the covariate vector, \(A\) is the binary treatment indicator, and \(Y(0),Y(1)\) are the binary potential responses under control and treatment. The prescribed functions \(e_\lambda,\mu_{0,\theta,\eta},\mu_{1,\theta,\eta}\) are the designated conditional-margin functions of this law.

\item Output the contrasts and finite pushforward measures
\[
\mathfrak C(h,\delta,a,b)
=
\bigl((\tau_\theta,\varpi_\theta)\bigr)_{\theta\in\{-1,+1\}},
\qquad
\varpi_\theta=
\sum_{\lambda,\eta\in\{-1,+1\}^{\mathcal I}}
w_\theta(\lambda,\eta)\,
\operatorname{Dirac}(P_{\theta,\lambda,\eta}),
\]
where \(\operatorname{Dirac}(P)\) is the unit point mass at the primitive law \(P\).
\end{enumerate}
\end{algorithmv}

The product inequality in \cref{obj:lem:hellinger-block-calculus} aggregates component distances. Its conditional identity integrates these distances over the common covariate distribution, subject to the stated measurable-density condition. The final identity accommodates the independent randomizer available to decisions in \cref{obj:def:risk}.

To translate mixture separation into expected absolute error, the analysis uses the fuzzy testing bound of \citet[Section 3, Lemma 1]{KennedyBalakrishnanRobinsWasserman2024}. \Cref{obj:lem:published-fuzzy-testing}, stated after the Hellinger calculus, fixes both its loss and its Hellinger normalization.

\begin{lemmav}{lem:hellinger-block-calculus}[Hellinger calculus for blocks]
For probability laws \(\mathbb B,\mathbb C\) on a common measurable space, use the normalization of \cref{obj:def:hellinger}:
\[
\mathsf h^2(\mathbb B,\mathbb C)
=\int\left(\sqrt{\frac{d\mathbb B}{d\xi}}
           -\sqrt{\frac{d\mathbb C}{d\xi}}\right)^2\,d\xi,
\qquad \xi=\mathbb B+\mathbb C.
\]
The following statements hold on arbitrary measurable spaces.

For any finite family of pairs of probability laws \(\mathbb B_{u_*},\mathbb C_{u_*}\), each pair on its respective common measurable space,
\[
\mathsf h^2\left(\bigotimes_{u_*}\mathbb B_{u_*},
                \bigotimes_{u_*}\mathbb C_{u_*}\right)
\le \sum_{u_*}\mathsf h^2(\mathbb B_{u_*},\mathbb C_{u_*}).
\]

Let \(\nu\) be a probability law on the conditioning space, and let \(\mathbb B_z,\mathbb C_z\) be probability kernels into a common measurable space. Suppose there exists a common s-finite reference kernel such that both kernels are represented, at every \(z\), by nonnegative extended-real densities with respect to that reference kernel, with both densities jointly measurable in the conditioning and outcome coordinates. Then
\[
\mathsf h^2\bigl(\nu(dz)\mathbb B_z(dw),
                \nu(dz)\mathbb C_z(dw)\bigr)
=\int \mathsf h^2(\mathbb B_z,\mathbb C_z)\,d\nu(z).
\]

For any probability laws \(\mathbb B,\mathbb C\) on a common measurable space and any probability law \(\nu\) on another measurable space,
\[
\mathsf h^2(\mathbb B\otimes\nu,\mathbb C\otimes\nu)
=\mathsf h^2(\mathbb B,\mathbb C).
\]
\end{lemmav}

For the original-record lower bound, the testing comparison concerns the law of the entire fixed-size dataset. At this level, the observation space contains both record blocks, and the testing statement can be used with product count one. This formulation preserves the unequal record counts in the experiment. The certificate also records an exact equality of the auxiliary mixture laws; the relevant marginal is the treatment-record marginal of \cref{obj:synth_4}.

\begin{lemmav}{lem:published-fuzzy-testing}[Absolute-loss fuzzy testing bound]
Let \((\mathbb B_z)\) and \((\mathbb C_z)\) be probability kernels from a common measurable index space to a common measurable observation space, and let \(\omega\) be a probability prior on the index space. Let \(v_0\ge1\) be an integer, let \(\mathcal M\) be a model of measures containing \(\mathbb B_z\) and \(\mathbb C_z\) for every \(z\), and let \(\Psi\) be a real-valued functional on measures. Use the squared Hellinger distance of \cref{obj:def:hellinger}, with normalization
\[
\mathsf h^2(\mathbb B,\mathbb C)
=
\int
\left(
\sqrt{\frac{d\mathbb B}{d\xi}}
-
\sqrt{\frac{d\mathbb C}{d\xi}}
\right)^2\,d\xi,
\qquad
\xi=\mathbb B+\mathbb C.
\]
Suppose the following conditions hold:
\begin{itemize}
\item \textbf{(Constants.)} The separation bound \(s_0\) and distance bound \(u_0\) satisfy \(s_0>0\) and \(0\le u_0<2\).
\item \textbf{(Separation.)} For every pair of indices \(z,z'\),
\[
\Psi(\mathbb B_z)-\Psi(\mathbb C_{z'})\ge s_0.
\]
\item \textbf{(Mixture distance.)} The product-law mixtures satisfy
\[
\mathsf h^2\left(
\int \mathbb B_z^{\otimes v_0}\,d\omega(z),
\int \mathbb C_z^{\otimes v_0}\,d\omega(z)
\right)\le u_0.
\]
\end{itemize}
Then every measurable real-valued decision \(\widehat\Psi\) on the \(v_0\)-observation sample space satisfies
\[
\sup_{P\in\mathcal M}
\int
\left|\widehat\Psi(w)-\Psi(P)\right|
\,dP^{\otimes v_0}(w)
\ge
\frac{s_0}{4}
\left\{1-\sqrt{u_0(1-u_0/4)}\right\},
\]
where the nonnegative loss integral is interpreted in the extended sense. This is the absolute-loss specialization of \citep[Section 3, Lemma 1]{KennedyBalakrishnanRobinsWasserman2024}.
\end{lemmav}

The marginal in \cref{obj:synth_4} is the sampling law of each outcome-free record. The next result establishes model membership, target separation, and the information bound for the prescribed smooth-frame construction. Its occupancy condition restricts the expected number of treatment records at the fine spatial scale.

\begin{theoremv}{thm:marked-component-certificate}[Marked component mixture certificate]
Fix the primitive class \(\mathcal M\) of \cref{obj:def:primitive-class}, with public parameters satisfying
\[
d\in\mathbb N,\qquad d\ge1,\qquad
\alpha,\beta\in(0,1],\qquad
\gamma\in[1,3],\qquad
L>\tfrac12,\qquad
\varepsilon\in(0,\tfrac12).
\]
There exist constants \(c\in(0,1]\), \(c_0>0\), and \(C>0\), depending only on these public parameters, such that the following holds for every choice satisfying:
\begin{itemize}
\item \textbf{(Record counts.)} \(n,m\in\mathbb N\), \(n\ge2\), with total record count \(N=n+m\).
\item \textbf{(Scales.)} \(0<\delta\le h\le\tfrac12\).
\item \textbf{(Propensity amplitude.)} \(0<a\le c\delta^\alpha\).
\item \textbf{(Response amplitude.)} \(0<b\le c\delta^\beta\).
\item \textbf{(Effect calibration.)} \(ab\le ch^\gamma\).
\item \textbf{(Occupancy.)} \(N\delta^d\le c_0\).
\end{itemize}
The prescribed construction \(\mathfrak C\) of \cref{obj:def:marked-handle}, evaluated at \(d,h,\delta,a,b\), has the following properties for both hypothesis indices \(\theta\in\{-1,+1\}\).

Its weights on the finite sign arrays are nonnegative and normalized:
\[
w_\theta(\lambda,\eta)\ge0
\quad\text{for every }\lambda,\eta\in\{-1,+1\}^{\mathcal I},
\qquad
\sum_{\lambda,\eta\in\{-1,+1\}^{\mathcal I}}
w_\theta(\lambda,\eta)=1.
\]
Every constructed law \(P_{\theta,\lambda,\eta}\) belongs to \(\mathcal M\), and its canonical contrast agrees with the deterministic constructed contrast everywhere on the unit cube:
\[
\tau_{P_{\theta,\lambda,\eta}}(x)=\tau_\theta(x)
\qquad
\bigl(x\in[0,1]^d\bigr).
\]
Thus the finite prior \(\varpi_\theta\) is supported on \(\mathcal M\). The designated propensity functions, restricted to the unit cube, have the same distribution under \(\varpi_{+1}\) and \(\varpi_{-1}\).

At the interior center \(x_0=(1/2,\ldots,1/2)\), the contrasts satisfy
\[
\tau_\theta(x_0)=-2\theta ab,
\qquad
\tau_{-1}(x_0)-\tau_{+1}(x_0)=4ab,
\qquad
cab\le4ab.
\]
For the fixed-size original-record mixture laws \(M_{+1}\) and \(M_{-1}\), the unhalved squared Hellinger distance of \cref{obj:def:hellinger} satisfies
\[
H^2(M_{+1},M_{-1})
\le CnNh^d\delta^da^2b^2.
\]

For every integer \(k_{\mathrm{aux}}\ge0\), the auxiliary-record mixtures agree:
\[
\int P_{XA}^{\otimes k_{\mathrm{aux}}}\,d\varpi_{+1}(P)
=
\int P_{XA}^{\otimes k_{\mathrm{aux}}}\,d\varpi_{-1}(P).
\]
Their explicit conditional auxiliary kernels also agree at every collection of \(k_{\mathrm{aux}}\) covariates in \([0,1]^d\). For either hypothesis and every \(x\in[0,1]^d\), the singleton labeled mixture conditional on \(X=x\) is uniform on the four values of \((A,Y)\in\{0,1\}^2\).

Conditional on every collection of all \(N\) covariates in the unit cube, form the graph whose vertices are all \(N\) records. Join two distinct vertices when both covariates lie in the localization cube \(C_h\) and their maximum-norm distance is at most \(2\delta\); records exterior to \(C_h\) are isolated vertices. The conditional mixed density equals the product of the component densities over all connected components of this graph, including the exterior singleton components. Every auxiliary-only component and every component consisting of a single labeled record has the same conditional density, and hence the same conditional law, under the two hypotheses.
\end{theoremv}

The information bound in \cref{obj:thm:marked-component-certificate} reflects where the two hypotheses can be distinguished. Their common propensity distribution gives identical auxiliary mixtures for every auxiliary sample size, including equality conditional on the auxiliary covariates. Compensation in the response means also makes a singleton labeled mixture uniform under either hypothesis. Conditional component factorization therefore locates distinguishing information in components containing a labeled record and at least one additional nearby record.

The factor \(nN\) in the distance bound corresponds to record pairs involving a labeled member, while \(h^d\delta^d\) records the localization and proximity scales. The squared perturbation factor is \(a^2b^2\), and the target separation is \(4ab\). Under the stated sparse-occupancy gate \(N\delta^d\le c_0\), with occupancy constant \leanref{sym:c0}{\ensuremath{c_0}}, the certificate bounds the aggregate information by this pair-scale expression. These clauses give the component structure needed by \cref{obj:lem:hellinger-block-calculus} and the separation needed by the testing bound of \cref{obj:lem:published-fuzzy-testing}.

The next result packages the smoothness, contrast, and occupancy calibrations into a minimax lower bound. Its sample-size interval is the range in which the unequal-count rate is at least the supplied-propensity order.

\begin{lemmav}{lem:nuisance-frontier-converse}[Nuisance frontier lower bound]
Fix \(d\in\mathbb N\) and \(\alpha,\beta,\gamma,L,\varepsilon\in\mathbb R\), where \(\varepsilon\) is the overlap level. Define the nuisance smoothness sum \(S\), the smoothness threshold \(S_{\mathrm{crit}}\), the record-growth exponent \(q_*\), and the rate denominator \(\Delta\) by
\[
S=\alpha+\beta,\qquad
S_{\mathrm{crit}}=\frac{\gamma d}{2\gamma+d},\qquad
q_*=\frac{\gamma d}{(\alpha+\beta)(2\gamma+d)},\qquad
\Delta=2\gamma+d+\frac{\gamma d}{\alpha+\beta}.
\]
Suppose that:
\begin{itemize}
\item \textbf{(Dimension.)} \(d\ge1\).
\item \textbf{(Nuisance smoothness.)} \(0<\alpha\le1\) and \(0<\beta\le1\).
\item \textbf{(Effect smoothness.)} \(1\le\gamma\le3\).
\item \textbf{(Radius.)} \(L>1/2\).
\item \textbf{(Overlap level.)} \(0<\varepsilon<1/2\).
\item \textbf{(Frontier regime.)} \(S<S_{\mathrm{crit}}\).
\end{itemize}
There exists a constant \(c>0\), depending on the fixed public parameters \(d,\alpha,\beta,\gamma,L,\varepsilon\), such that for every \(n,m\in\mathbb N\) with \(n\ge2\), the total treatment-record count
\[
N=n+m
\]
satisfying \(N\le n^{q_*}\) obeys the following bound for the minimax absolute-error risk \(R(n,m)\) of \cref{obj:def:minimax}, evaluated at these fixed parameters:
\[
R(n,m)\ge c(nN)^{-\gamma/\Delta}.
\]
\end{lemmav}

Under the low-nuisance-smoothness condition in \cref{obj:lem:nuisance-frontier-converse}, the exponent \(q_*\) exceeds one. The stated interval consequently covers total record counts from the fully labeled experiment through the crossover to supplied-propensity precision. Within this interval, the lower bound captures the joint contribution of labeled and outcome-free records through their count product.

The complete minimax characterization in \cref{obj:thm:sharp-annotation-frontier} organizes the lower bounds into complementary ranges. When \(S<S_{\mathrm{crit}}\) and \(n\le N\le n^{q_*}\), \cref{obj:lem:nuisance-frontier-converse} supplies the unequal-count term. Supplying the propensity enlarges the information available to a decision, so the supplied-propensity risk in \cref{obj:def:oracle-risk} is a lower bound for the original-record minimax risk in \cref{obj:def:minimax}. For \(N\ge n^{q_*}\), the supplied-propensity lower bound in \cref{obj:thm:supplied-propensity} supplies the order \(n^{-\gamma/(2\gamma+d)}\), which is the larger term in the sharp rate. When \(S\ge S_{\mathrm{crit}}\), that same order is the larger term for every \(N\ge n\). Thus the supplied-propensity floor resolves the ranges complementary to the sparse-occupancy construction.

Together with \cref{obj:thm:rectangular-upper}, these lower bounds yield the matching characterization in \cref{obj:thm:sharp-annotation-frontier}. For the formal crosswalk, the rate and attaining decision are collected in the ordered pair of \cref{obj:def:frontier-handle} in \cref{sec:verification-note}.



The statistical interpretation and count consequences of this pair are stated once in \cref{obj:thm:sharp-annotation-frontier,obj:prop:annotation-threshold}.
